\documentclass[12pt, a4paper]{article}

% --- Pakiety podstawowe ---
\usepackage[utf8]{inputenc}
\usepackage[T1]{fontenc}
\usepackage[polish]{babel}
\usepackage{polski}
\usepackage{geometry}
\usepackage{graphicx}
\usepackage{float}
\usepackage{hyperref}
\usepackage{amsmath, amsfonts, amssymb}
\usepackage{listings}
\usepackage{xcolor}
\usepackage{caption}
\usepackage{subcaption}

% --- Konfiguracja marginesów ---
\geometry{
 a4paper,
 total={170mm,257mm},
 left=20mm,
 top=20mm,
}

% --- Konfiguracja wyświetlania kodu (opcjonalnie) ---
\definecolor{codegray}{rgb}{0.5,0.5,0.5}
\definecolor{backcolour}{rgb}{0.95,0.95,0.92}

\lstdefinestyle{mystyle}{
    backgroundcolor=\color{backcolour},   
    commentstyle=\color{codegray},
    keywordstyle=\color{blue},
    numberstyle=\tiny\color{codegray},
    stringstyle=\color{green!50!black},
    basicstyle=\ttfamily\footnotesize,
    breakatwhitespace=false,         
    breaklines=true,                 
    captionpos=b,                    
    keepspaces=true,                 
    numbers=left,                    
    numbersep=5pt,                  
    showspaces=false,                
    showstringspaces=false,
    showtabs=false,                  
    tabsize=2
}
\lstset{style=mystyle}

% --- Dane dokumentu ---
\title{\textbf{Symulacja rozprzestrzeniania się epidemii zombie}\\
\large Projekt z przedmiotu Modelowanie Matematyczne}
\author{Patryk Buchtyar, Rafał Maciończyk, Michał Moryc}
\date{\today}

\begin{document}

\maketitle

\begin{abstract}
    Niniejszy projekt przedstawia symulację wieloagentową modelującą interakcje między czterema grupami: ludźmi, zombie, medykami oraz żołnierzami. Celem pracy jest analiza dynamiki populacji w zamkniętym środowisku oraz ocena skuteczności algorytmów uczenia ze wzmocnieniem (Deep Q-Learning) w sterowaniu zachowaniem agentów. Porównano wyniki symulacji z klasycznymi modelami matematycznymi, takimi jak model Lotki-Volterry.
\end{abstract}

\tableofcontents
\newpage

\section{Wstęp}
W tym rozdziale należy opisać cel projektu. Dlaczego symulujemy zombie? (np.
jako metaforę agresywnej epidemii). Jakie są założenia środowiska (siatka, czas
dyskretny).

\section{Opis Modelu Symulacji}
Symulacja odbywa się w środowisku dyskretnym. Wyróżniamy cztery klasy agentów,
z których każda posiada unikalne mechaniki:

\begin{itemize}
    \item \textbf{Ludzie:} Ich celem jest przetrwanie. Posiadają zdolność reprodukcji z prawdopodobieństwem $P_{rep} = 0.002$ na krok. Istnieje również szansa ($20\%$), że nowo narodzony człowiek przejdzie mutację społeczną i stanie się Medykiem lub Żołnierzem.
    \item \textbf{Zombie:} Dążą do zarażenia każdego napotkanego agenta. Posiadają mechanizm głodu (starvation) – brak pożywienia przez określony czas skutkuje śmiercią agenta. Infekcja posiada czas odnowienia (cooldown).
    \item \textbf{Żołnierze:} Ich zadaniem jest eliminacja zombie. Mechanika ataku jest ograniczona czasowo (cooldown).
    \item \textbf{Medycy:} Agentzy specjalni, których celem jest leczenie zarażonych. Poruszają się, wykorzystując deterministyczny algorytm ścieżki (opisany w sekcji 3).
\end{itemize}

\section{Sterowanie i Algorytmy}

W symulacji zastosowano podejście hybrydowe. Medycy wykorzystują klasyczne
algorytmy przeszukiwania grafu, natomiast pozostałe klasy (Ludzie, Zombie,
Żołnierze) są sterowane przez głębokie sieci neuronowe uczone metodą
\textit{Deep Q-Network} (DQN).

\subsection{Algorytm A* (Medycy)}
Medycy nie wymagają procesu uczenia, ponieważ ich cel jest jasno zdefiniowany i
statyczny. Do znajdowania najkrótszej ścieżki do najbliższego zarażonego agenta
wykorzystano algorytm A* (A-star). Funkcja kosztu $f(n)$ zdefiniowana jest
jako: $$ f(n) = g(n) + h(n) $$ gdzie $g(n)$ to koszt dotarcia do węzła $n$, a
$h(n)$ to heurystyka (w tym przypadku odległość Manhattan lub Euklidesowa).

\subsection{Głębokie Uczenie ze Wzmocnieniem (DQN)}

[Image of convolutional neural network architecture diagram]

Dla agentów sterowanych sztuczną inteligencją zastosowano algorytm Deep
Q-Learning. Agent uczy się aproksymować funkcję wartości $Q(s, a)$, która
określa przewidywaną sumę nagród dla podjęcia akcji $a$ w stanie $s$.

\subsubsection{Architektura Sieci Neuronowej (Sensor Fusion)}
Model decyzyjny oparty jest na architekturze typu \textit{Sensor Fusion}, która
łączy dane wizualne (lokalne otoczenie agenta) z danymi wektorowymi (np.
względna pozycja celu). Sieć zaimplementowana w bibliotece \texttt{PyTorch}
składa się z dwóch gałęzi:

\begin{enumerate}
    \item \textbf{Gałąź Konwolucyjna (CNN):} Przetwarza lokalną mapę otoczenia agenta.
          \begin{itemize}
              \item \textbf{Wejście:} Tensor o wymiarach $(4, 11, 11)$, co odpowiada 4 kanałom (np. klatki czasowe lub typy obiektów) na siatce $11 \times 11$.
              \item \textbf{Warstwa Conv1:} 32 filtry, kernel $3 \times 3$, padding 1, aktywacja ReLU.
              \item \textbf{Warstwa Conv2:} 64 filtry, kernel $3 \times 3$, padding 1, aktywacja ReLU.
              \item \textbf{Wyjście:} Spłaszczony wektor cech wizualnych o wymiarze $64 \times 11 \times 11$.
          \end{itemize}

    \item \textbf{Gałąź Wektorowa (MLP):} Przetwarza proste dane numeryczne.
          \begin{itemize}
              \item \textbf{Wejście:} Wektor wielkości 2 (np. $\Delta x, \Delta y$ do najbliższego celu/zagrożenia).
              \item \textbf{Warstwa liniowa:} Transformacja do 32 cech z aktywacją ReLU.
          \end{itemize}

    \item \textbf{Fuzja i Decyzja:}
          Cechy z obu gałęzi są konkatenowane, tworząc wspólny wektor reprezentacji stanu. Następnie przechodzą przez warstwy gęste (Fully Connected):
          \begin{itemize}
              \item Warstwa ukryta: 512 neuronów.
              \item Warstwa wyjściowa: Liczba neuronów równa liczbie możliwych akcji (5),
                    zwracająca wartości $Q$ dla każdej akcji.
          \end{itemize}
\end{enumerate}

\subsubsection{Proces Uczenia}
Uczenie odbywa się przy użyciu bufora doświadczeń (\textit{Replay Buffer}) o
pojemności 10,000 próbek. Funkcja straty (Loss) to Błąd Średniokwadratowy (MSE)
pomiędzy przewidywaną wartością Q, a wartością docelową obliczaną z równania
Bellmana:

$$ L(\theta) = \mathbb{E} \left[ \left( r + \gamma \max_{a'} Q(s', a'; \theta^{-}) - Q(s, a; \theta) \right)^2 \right] $$

Gdzie:
\begin{itemize}
    \item $\gamma = 0.99$ (współczynnik dyskontowania),
    \item $\theta$ to wagi sieci głównej (Policy Net),
    \item $\theta^{-}$ to wagi sieci docelowej (Target Net), aktualizowane co 1000 kroków,
    \item Strategia eksploracji: $\epsilon$-greedy z zanikiem (decay) od $\epsilon=1.0$
          do $\epsilon_{min}=0.05$.
\end{itemize}

\section{Modele Matematyczne Dynamiki Populacji}
W celu weryfikacji realizmu symulacji, uzyskane wyniki przyrostu populacji
porównano z teoretycznymi modelami matematycznymi.

\subsection{Model Lotki-Volterry}
Relacja między Żołnierzami/Ludźmi (ofiary) a Zombie (drapieżniki) może być
opisana układem równań różniczkowych pierwszego rzędu:

$$
    \begin{cases}
        \frac{dx}{dt} = \alpha x - \beta xy \\
        \frac{dy}{dt} = \delta xy - \gamma y
    \end{cases}
$$

gdzie:
\begin{itemize}
    \item $x$ -- liczba ludzi/żołnierzy,
    \item $y$ -- liczba zombie,
    \item $\alpha, \beta, \delta, \gamma$ -- parametry modelu estymowane na podstawie danych z symulacji.
\end{itemize}

\subsection{Logistyczny Model Wzrostu}
Dla populacji ludzi, ograniczonej zasobami (miejscem na planszy), stosuje się
model logistyczny: $$ \frac{dP}{dt} = rP \left(1 - \frac{P}{K}\right) $$ gdzie
$K$ to pojemność środowiska (carrying capacity).

\section{Wyniki Symulacji}
 (Tutaj wstawisz wykresy wygenerowane z Matplotlib, np. liczebność populacji w czasie).

\section{Wnioski}
Podsumowanie projektu. Czy sieć neuronowa nauczyła się skutecznych strategii?
Czy populacja zombie zdominowała świat, czy została opanowana?

\end{document}